In this note we prove a simple lower bound for the numbers
that appears to be asymptotically exact as $g\to\infty$. To
begin with, we recall some basic facts about the numbers
$\langle\tau_{d_1}\cdots\tau_{d_n}\rangle$ (Witten's
correlators). They are uniquely defined by the initial data
\[
\big\langle \tau_0^3 \big\rangle = 1,\qquad
\langle \tau_{1} \rangle = \frac{1}{24}
\]
via the recursive relations known as Virasoro constraints that we present below.

\textbf{Virasoro constraints} (in Dijkgraaf--Verlinde--Verlinde form~\cite{Dij, DVV}):
\begin{gather}
\langle\tau_{k+1}\tau_{d_1}\cdots\tau_{d_n}\rangle_g=\frac{1}{(2k+3)!!}\Bigg[\sum_{j=1}^n
\frac{(2k+2d_j+1)!!}{(2d_j-1)!!}\langle\tau_{d_1}\cdots
\tau_{d_{j}+k}\cdots\tau_{d_n}\rangle_g\nonumber\\
\hphantom{\langle\tau_{k+1}\tau_{d_1}\cdots\tau_{d_n}\rangle_g=}{} +\frac{1}{2}\sum_{\substack{r+s=k-1\\r,s\ge 0}}
(2r+1)!!(2s+1)!!\langle\tau_r\tau_s\tau_{d_1}\cdots\tau_{d_n}\rangle_{g-1}\nonumber\\
\hphantom{\langle\tau_{k+1}\tau_{d_1}\cdots\tau_{d_n}\rangle_g=}{} +\frac{1}{2}\sum_{\substack{r+s=k-1\\r,s\ge 0}} (2r+1)!!(2s+1)!!\nonumber\\
\hphantom{\langle\tau_{k+1}\tau_{d_1}\cdots\tau_{d_n}\rangle_g=}{}\times
\sum_{\{1,\dots,n\}=I\coprod J}\bigg\langle\tau_r\prod_{i\in
I}\tau_{d_i}\bigg\rangle_{g'}\bigg\langle\tau_s\prod_{i\in
J}\tau_{d_i}\bigg\rangle_{g-g'}\Bigg].\label{eq:virasoro}
\end{gather}

For $k=-1$ and $k=0$ the above relations have particularly simple form.

\textbf{String equation ($\boldsymbol{k=-1}$):}
\begin{gather}\label{eq:string}
\langle \tau_0\tau_{d_1} \cdots \tau_{d_n}\rangle_{g,n+1}
=\langle \tau_{d_1-1} \cdots \tau_{d_n}\rangle_{g,n}
+\dots
+\langle \tau_{d_1} \cdots \tau_{d_n-1}\rangle_{g,n} .
\end{gather}

\textbf{Dilaton equation ($\boldsymbol{k=0}$):}
\begin{gather*}%\label{eq:dilaton}
\langle \tau_1\tau_{d_1} \cdots \tau_{d_n}\rangle_{g,n+1}
=(2g-2+n)\langle \tau_{d_1} \cdots \tau_{d_n}\rangle_{g,n} .
\end{gather*}

For any partition $\boldsymbol{d}$ of $3g-3+n$
into a sum of $n$ nonnegative integers define
$\epsilon(\boldsymbol{d})$ through the following equation:
\begin{gather*}%\label{eq:ansatz}
\big\langle \psi_1^{d_1} \cdots \psi_n^{d_n}\big\rangle_{g,n}
=
\frac{(6g-5+2n)!!}{(2d_1+1)!!\cdots(2d_n+1)!!}
\cdot\frac{1}{g!\cdot 24^g} \cdot\big(1+\epsilon(\boldsymbol{d})\big) .
\end{gather*}

We denote by $\Pi(m,n)$ the
set of ordered partitions of an integer $m$ into a sum of $n$ nonnegative integers.
